Angular momentum from a circular Keplerian accretion disk around a young star
of mass $M$ is removed by stellar winds. Wind particles slide frictionlessly
along magnetic field lines corotating with the disk. Consider a wind particle
of mass $m$ at radius $r$ and height $z$ moving along a magnetic field line
that corotates with the disk at radius $r_0$ and is inclined at an angle
$\theta_0$.

\ioaafig{2026_TH_T03-f1.pdf}

\begin{description}
\item[(a)] (2 points) In the frame corotating with the particle, derive
  $\Phi_{\mathrm{eff}}(r, z)$, defined as the potential energy from the
  star's gravity and centrifugal force.
\item[(b)] (5 points) Introducing the variable $x = \frac{r - r_0}{r_0}$,
  write $\Phi_{\mathrm{eff}}$ in the vicinity of $r_0$ to second-order
  approximation in $x$. Find the inclination angle $\theta_0$ so that the
  particle can be pulled away from the disk. Hint:
  $(1 + u)^{-1/2} \approx 1 - \frac{u}{2} + \frac{3}{8}u^2$ with
  $|u| \ll 1$.
\item[(c)] (3 points) At radius $r = r_{\mathrm{A}}$, particles begin to
  decouple from magnetic field. The wind launched from $r_0$ carries
  mass-loss rate $\Delta M_{\mathrm{w}}$ while the accretion rate is
  $\Delta M_{\mathrm{acc}}$. If all accretion angular momentum is removed by
  the wind, show that
  $\frac{\Delta M_{\mathrm{acc}}}{\Delta M_{\mathrm{w}}}
   = \left(\frac{r_{\mathrm{A}}}{r_0}\right)^2$.
\end{description}
